% Realization Calculus / professional book design.
% Fonts are system dependencies, not bundled assets.
\definecolor{RCPlum}{HTML}{39263F}
\definecolor{RCMulberry}{HTML}{75465F}
\definecolor{RCGold}{HTML}{A78A50}
\definecolor{RCInk}{HTML}{29252B}
\definecolor{RCMuted}{HTML}{68616C}
\definecolor{RCPale}{HTML}{DDD6DF}
\colorlet{ink}{RCInk}
\colorlet{accent}{RCMulberry}
\colorlet{navy}{RCPlum}
\colorlet{muted}{RCMuted}
\colorlet{warm}{RCGold}
\color{RCInk}
\setlength{\parindent}{0pt}
\setlength{\parskip}{6pt plus 1pt minus .5pt}
\linespread{1.10}
\setlength{\emergencystretch}{2.5em}
\setlength{\abovedisplayskip}{10pt plus 2pt minus 2pt}
\setlength{\belowdisplayskip}{10pt plus 2pt minus 2pt}
\setlength{\abovedisplayshortskip}{6pt plus 2pt}
\setlength{\belowdisplayshortskip}{8pt plus 2pt minus 2pt}
\setlength{\jot}{4pt}
\raggedbottom
\clubpenalty=10000
\widowpenalty=10000
\displaywidowpenalty=10000
\interfootnotelinepenalty=10000
\setcounter{tocdepth}{0}
\setcounter{secnumdepth}{2}
\newif\ifRCAppendix
\newcommand{\SmallLabel}[1]{{\sffamily\addfontfeatures{LetterSpace=13}\fontsize{8}{11}\selectfont\color{RCMulberry}\MakeUppercase{#1}}}
\newcommand{\GoldRule}{\par\vspace{10pt}{\color{RCGold}\rule{.56in}{.65pt}}\par\vspace{13pt}}
\newcommand{\CurrentChapter}{SPECIAL FUNCTIONS}
\pagestyle{fancy}
\fancyhf{}
\fancyhead[L]{\sffamily\addfontfeatures{LetterSpace=5}\fontsize{7.1}{10}\selectfont\color{RCMuted}REALIZATION CALCULUS}
\fancyhead[R]{\sffamily\fontsize{7.1}{10}\selectfont\color{RCMuted}\MakeUppercase{\CurrentChapter}}
\fancyfoot[R]{\rmfamily\fontsize{10}{12}\selectfont\color{RCInk}\thepage}
\renewcommand{\headrulewidth}{0pt}
\renewcommand{\footrulewidth}{0pt}
\fancypagestyle{plain}{\fancyhf{}\fancyfoot[R]{\rmfamily\fontsize{10}{12}\selectfont\color{RCInk}\thepage}\renewcommand{\headrulewidth}{0pt}}
\renewcommand{\chaptermark}[1]{\gdef\CurrentChapter{#1}}
\titleformat{\chapter}[display]
 {\normalfont\raggedright\hyphenpenalty=10000\exhyphenpenalty=10000\color{RCPlum}}
 {\SmallLabel{\ifRCAppendix Appendix \thechapter\else Chapter \Roman{chapter}\fi}}
 {20pt}{\fontsize{31}{37}\selectfont}[\GoldRule]
\titlespacing*{\chapter}{0pt}{19pt}{12pt}
\titleformat{\section}
 {\normalfont\raggedright\hyphenpenalty=10000\exhyphenpenalty=10000\fontsize{18}{23}\selectfont\color{RCPlum}}
 {\textcolor{RCMulberry}{\thesection}}{.65em}{}
\titlespacing*{\section}{0pt}{23pt plus 4pt minus 2pt}{9pt}
\titleformat{\subsection}
 {\normalfont\raggedright\hyphenpenalty=10000\fontsize{12.4}{16}\selectfont\bfseries\color{RCInk}}
 {\thesubsection}{.6em}{}
\titlespacing*{\subsection}{0pt}{15pt plus 3pt minus 2pt}{5pt}
\titleformat{\subsubsection}{\normalfont\bfseries\color{RCInk}}{}{0pt}{}
\titlespacing*{\subsubsection}{0pt}{12pt}{4pt}
\hypersetup{colorlinks=true,linkcolor=RCPlum,citecolor=RCPlum,urlcolor=RCMulberry,
 pdftitle={Special Functions: Function Families and Their Structure},
 pdfauthor={Komorebi P.},pdfsubject={Realization Calculus / Special Functions},
 bookmarksopen=true,bookmarksopenlevel=0,bookmarksnumbered=true,bookmarksdepth=2,pdfdisplaydoctitle=true}
\captionsetup{hypcap=false,font=small,labelfont={bf},labelsep=period,justification=justified,singlelinecheck=false,skip=8pt}
\setlength{\textfloatsep}{16pt plus 2pt minus 2pt}
\setlength{\intextsep}{15pt plus 2pt minus 2pt}
\setlist[itemize]{leftmargin=1.4em,itemsep=3pt,topsep=5pt}
\setlist[enumerate]{leftmargin=1.7em,itemsep=3pt,topsep=5pt}
\setlist[description]{style=nextline,leftmargin=0pt,labelsep=0pt,font=\normalfont\bfseries\color{RCPlum}}
\renewcommand{\arraystretch}{1.22}
\setlength{\tabcolsep}{5pt}
\makeatletter
\newcommand{\pageheading}[2]{\if@mainmatter\section{#2}\else\section*{#2}\fi}
\makeatother
\newcommand{\claim}[1]{\par\Needspace{4\baselineskip}\vspace{7pt}{\normalfont\bfseries\color{RCPlum}#1}\par\nobreak\vspace{1pt}}
\newcommand{\fig}[2]{\par\addvspace{11pt}\noindent\begin{minipage}{\linewidth}\centering\includegraphics[width=\linewidth]{figures/#1.pdf}\captionof{figure}{#2}\label{fig:#1}\end{minipage}\par\addvspace{11pt}}
\newcommand{\smallfig}[3][.82]{\par\addvspace{11pt}\noindent\begin{minipage}{\linewidth}\centering\includegraphics[width=#1\linewidth]{figures/#2.pdf}\captionof{figure}{#3}\label{fig:#2}\end{minipage}\par\addvspace{11pt}}
\newcommand{\vfig}[3][.96]{\par\addvspace{11pt}\noindent\begin{minipage}{\linewidth}\centering\includegraphics[width=#1\linewidth]{figures_v11/#2.pdf}\captionof{figure}{#3}\label{fig:#2}\end{minipage}\par\addvspace{11pt}}
\newcommand{\FrontChapter}[1]{\cleardoublepage\chapter*{#1}\phantomsection\addcontentsline{toc}{chapter}{#1}\gdef\CurrentChapter{#1}}
\newcommand{\MonographCover}{%
 \hypersetup{pageanchor=false}\thispagestyle{empty}\begingroup
 \begin{tikzpicture}[remember picture,overlay]
 \fill[RCPlum] (current page.south west) rectangle (current page.north east);
 \node[anchor=north west,text=white,font=\sffamily\addfontfeatures{LetterSpace=16}\fontsize{9}{12}\selectfont] at ([xshift=.9in,yshift=-1in]current page.north west) {REALIZATION CALCULUS};
 \draw[RCGold,line width=.8pt] ([xshift=.92in,yshift=-1.52in]current page.north west) -- ++(.56in,0);
 \node[anchor=north west,text=white,align=left,inner sep=0pt,font=\rmfamily\fontsize{59}{66}\selectfont] at ([xshift=.88in,yshift=-1.92in]current page.north west) {Special\\Functions};
 \node[anchor=north west,text=white,align=left,inner sep=0pt,font=\rmfamily\fontsize{23}{29}\selectfont] at ([xshift=.92in,yshift=-4.32in]current page.north west) {Function families\\and their structure};
 \node[anchor=north west,text=RCPale,align=left,inner sep=0pt,font=\rmfamily\fontsize{13}{19}\selectfont] at ([xshift=.92in,yshift=-5.53in]current page.north west) {Coupled growth, arithmetic depth\\and finite resolution};
 \draw[white!25!RCPlum,line width=.45pt] ([xshift=.92in,yshift=1.29in]current page.south west) -- ([xshift=-.88in,yshift=1.29in]current page.south east);
 \node[anchor=south west,text=white,inner sep=0pt,font=\rmfamily\fontsize{14}{18}\selectfont] at ([xshift=.92in,yshift=.82in]current page.south west) {Komorebi P.};
 \end{tikzpicture}\null\endgroup\clearpage\hypersetup{pageanchor=true}
}
\newcommand{\MonographTitleLeaf}{%
 \thispagestyle{plain}\vspace*{13pt}
 \SmallLabel{Realization Calculus}\par
 \vspace{59pt}
 {\raggedright\fontsize{44}{51}\selectfont\color{RCPlum}Special\\Functions\par}
 \vspace{25pt}
 {\raggedright\fontsize{23}{29}\selectfont\color{RCPlum}Function families\\and their structure\par}
 \GoldRule
 {\fontsize{13}{19}\selectfont\color{RCMuted}Coupled growth, arithmetic depth\\and finite resolution\par}
 \vspace{76pt}
 {\fontsize{14}{19}\selectfont Komorebi P.\par}
 \vspace{22pt}
 {\fontsize{10.5}{15}\selectfont\color{RCMuted}A mathematical monograph\par}
 \clearpage
}
\newcommand{\MonographPreface}{%
 \FrontChapter{Preface}
 \section*{Structure beyond a finite window}
 The theory of special functions begins with explicit definitions, but the information those definitions carry need not be visible in a small sample of values. A contraction in a differential equation can determine the growth of an entire function. A prime omitted from an arithmetic sum can change its exact parameter signature. A finite collection of positive moment matrices can conceal a negative mass in the underlying measure.

 This volume studies these distinctions through five families. Scalar nonlinear feedback leads to entire functions with explicit growth laws. Arithmetic depth leads to positive kernels and an exact inverse problem. Prime sieves give a setting in which positivity and finite detection can be separated. Zero orbits provide spectral coordinates. Coupled contraction functions bring local coefficients, global growth and finite observation into one system.

 The first seven chapters develop the constructions, their principal arguments and worked examples. Ten appendices supply the longer proofs and computational details. The evidence register, glossary and references follow at the back. Figures accompany the mathematics they explain, rather than forming a separate gallery.

 Each result is read with its own hypotheses. Classical inputs are identified where they enter. Exact rational certificates, numerical illustrations and analytic proofs remain distinct, and unresolved questions retain their stated scope. A change of representation does not, by itself, establish historical priority.

 \vspace{7pt}
 The recurring question is precise: what can be recovered from finite information, and what additional resolution is needed to distinguish the underlying structure?
 \clearpage
}
\newcommand{\MonographPart}[3]{%
 \clearpage\refstepcounter{part}\phantomsection\addcontentsline{toc}{part}{\texorpdfstring{Part #1\quad #2}{Part #1: #2}}\thispagestyle{plain}
 \gdef\CurrentChapter{SPECIAL FUNCTIONS}
 \vspace*{91pt}\SmallLabel{Part #1}\par\vspace{22pt}
 {\raggedright\hyphenpenalty=10000\fontsize{33}{40}\selectfont\color{RCPlum}#2\par}
 \vspace{13pt}\GoldRule
 \begin{minipage}{.85\linewidth}\fontsize{12.2}{18}\selectfont\color{RCMuted}#3\end{minipage}
 \vfill\SmallLabel{Special Functions}\par\clearpage
}
% Contents: restrained typography, generous hierarchy and live links.
\makeatletter
\renewcommand{\l@part}[2]{\addpenalty{-\@highpenalty}\addvspace{15pt}\begingroup\parindent\z@\rightskip\@pnumwidth\parfillskip-\@pnumwidth\normalfont\bfseries\color{RCPlum}#1\nobreak\hfil\nobreak\hb@xt@\@pnumwidth{\hss #2}\par\endgroup\nobreak}
\renewcommand{\l@chapter}[2]{\addpenalty{-\@highpenalty}\vskip5pt\begingroup\parindent\z@\rightskip\@pnumwidth\parfillskip-\@pnumwidth\normalfont\leftskip1.5em\hskip-\leftskip\def\numberline##1{\hb@xt@2.2em{\color{RCMulberry}##1\hfil}}#1\nobreak\hfil\nobreak\hb@xt@\@pnumwidth{\hss #2}\par\endgroup}
\makeatother
\newcolumntype{P}[1]{>{\raggedright\arraybackslash}p{#1}}
